\chapter{Quên Bút, Quên Thước}
\chapterintrobi{Đi học đủ người nhưng thiếu đồ nghề}{Forget the pen and ruler}{Thói quen này khiến tiết học nào cũng bắt đầu bằng một vòng mượn khẩn cấp.}

\begin{lucbatbox}{Lục bát: Cặp to nhưng vẫn thiếu}
Cặp em tưởng đủ trăm phần\
Mở ra mới biết thiếu dần đồ quen\
Bút đâu, thước ở nơi trên\
Hay là ở lại cùng đêm hôm nào\
Bạn bên vừa học vừa trao\
Thành ra kiến thức đi vào nhờ vay
\end{lucbatbox}

\begin{tutuyetbox}{Thất ngôn tứ tuyệt: Mượn đồ đầu giờ}
Chuông vừa điểm đã ngó quanh bàn\
Bạn ơi cho mượn chiếc bút sang\
Nếu em nhớ đồ như nhớ chuyện nhỏ\
Chắc đời đi học đỡ hoang mang
\end{tutuyetbox}

\begin{ghichu}
Quên bút quên thước không phải bi kịch lớn, nhưng lặp lại mãi thì nó cho thấy việc chuẩn bị đang bị xem hơi nhẹ.
\end{ghichu}

\begin{englishnote}
\textbf{borrow a pen}: mượn bút.\par
\textbf{school supplies}: đồ dùng học tập.\par
\textbf{I forgot my ruler again.}: Em lại quên thước rồi.
\end{englishnote}